%     - Měření SQL dotazů: data z `rust-lang/rust`, jak dlouho trvá konkrétních 5 dotazů? Pomáhají indexy? `EXPLAIN ANALYZE` výstup.

\section{Měření výkonu analytických dotazů}
\label{sec:evaluace:sql}

V~rámci evaluace bylo měřeno, jaký vliv mají databázové indexy na dobu vykonávání analytických dotazů. Databáze obsahovala data repozitáře \texttt{rust-lang/rust}: přibližně 200\,000 záznamů issues a~\PR, 564\,000 záznamů v~tabulce \code{issue\_event\_history}, 355\,000 v~\code{issue\_labels\_history} a~525\,000 v~\code{file\_activity}. Měření bylo provedeno pomocí \code{EXPLAIN ANALYZE} v~PostgreSQL.

\subsection{Testované dotazy}
\label{subsec:evaluace:sql:dotazy}

Byly testovány čtyři dotazy (S1--S4) pokrývající typické analytické scénáře. Dotazy jsou uvedeny v~\customref{Příloze}{chap:appendix:sql}. Stručný popis každého dotazu:

\begin{description}
  \item[S1] Počet \PR ve stavu \code{S-waiting-on-review} k~zadanému datu. Dotaz využívá dvě \CTE s~\code{DISTINCT ON}: jednu pro poslední label-event a~druhou pro poslední stavovou změnu, následované \code{LEFT JOIN} a~filtrací.
  \item[S2] Kumulativní počet sloučených \PR k~zadanému datu pomocí \code{COUNT(DISTINCT issue)} s~filtrem \code{event = 'merged'} nad tabulkou \code{issue\_event\_history}.
  \item[S3] Soubory upravené přispěvateli z~daného týmu v~zadaném časovém rozsahu, agregace přes \code{file\_activity} a~\code{contributors\_teams}.
  \item[S4] Denní vývoj počtu \PR v~daném stavu v~časovém rozsahu. Nejsložitější dotaz využívá delta-event přístup: nejprve sestaví intervaly otevřenosti a~aktivity štítku a~vypočítá jejich průnik (operátory \code{\&\&} a~\code{*}). Intervaly jsou následně rozloženy na delta-události (+1/$-1$) a~nad nimi je spočtán kumulativní součet. Výsledná složitost je $O(\text{events} + \text{days})$ místo naivního range-joinu.
\end{description}

\subsection{Vliv indexů}
\label{subsec:evaluace:sql:indexy}

Pro optimalizaci byly vytvořeny částečné a~pokrývající indexy. Úplné definice indexů jsou uvedeny v~\customref{Příloze}{chap:appendix:sql}. Stručně:

\begin{itemize}
  \item \code{idx\_pr\_event\_hist\_state\_changes} -- částečný index na \code{issue\_event\_history} s~podmínkou \code{WHERE is\_pr = true AND event IN ('closed','merged','reopened')} a~klauzulí \code{INCLUDE (event)}.
  \item \code{idx\_pr\_labels\_repo\_label\_issue\_ts} -- pokrývající index na \code{issue\_labels\_history} se sloupci \code{(repository, label, issue, timestamp DESC) INCLUDE (action)}.
  \item \code{idx\_file\_activity\_repo\_contrib\_ts} -- pokrývající index na \code{file\_activity} se sloupci \code{(repository, contributor\_id, timestamp DESC) INCLUDE (file\_path, issue)}.
\end{itemize}

Částečné indexy s~podmínkou \code{WHERE is\_pr = true} redukovaly velikost indexu přibližně na polovinu, protože tabulka obsahuje záznamy jak pro issues, tak pro \PR. Pokrývající indexy s~klauzulí \code{INCLUDE} umožnily Index-Only Scan bez přístupu k~datovým stránkám (\code{Heap Fetches: 0}).

\subsection{Výsledky měření}
\label{subsec:evaluace:sql:vysledky}

Výsledky měření jsou shrnuty v~\customref{Tabulce}{tab:evaluace:sql}, plné výstupy \code{EXPLAIN (ANALYZE, BUFFERS)} jsou k~dispozici \customref{v~příloze}{sec:appendix:sql:explain}.

Nejvýraznější přínos měly indexy u~S1 a~S2. Bez indexů musel plánovač provést \code{Seq Scan} celé tabulky a~výsledky seřadit na disku (\code{external merge Disk}) \listingref{lst:appendix:explain:s1:no_idx}, \listingref{lst:appendix:explain:s2:no_idx}.
Se~zavedením pokrývajících indexů přešel plán na \code{Index Only Scan} \listingref{lst:appendix:explain:s1:idx}, \listingref{lst:appendix:explain:s2:idx} a~třídění na disku zcela odpadlo.
Výsledkem je 3,7$\times$, resp. 7,0$\times$ kratší čas vykonávání.

U~S3 byl \code{Parallel Seq Scan} přes celou tabulku \code{file\_activity} nahrazen
\code{Nested Loop} \listingref{lst:appendix:explain:s3:no_idx}, \listingref{lst:appendix:explain:s3:idx}. Ten pro každého člena týmu načítá záznamy přímo z~pokrývajícího indexu,
tudíž odpadá přístup k~datovým stránkám (\code{Heap Fetches: 0}).
Zrychlení je v~porovnání s~S1 a~S2 menší (1,2$\times$), protože agregace a~třídění výsledků tvoří větší podíl celkového času.

S4 zaznamenal zlepšení 1,5$\times$ \listingref{lst:appendix:explain:s4:no_idx}, \listingref{lst:appendix:explain:s4:idx}. Hlavní náklady tvoří \code{Merge Join} při výpočtu průniku intervalů (\code{in\_state\_periods}), který indexovat nelze. Přesto indexy přinesly úsporu: sekvenční průchody oběma tabulkami byly nahrazeny \code{Index Only Scan} \\a~\code{Bitmap Heap Scan} a~řazení pro okenní funkce přešlo z~external merge na quicksort v~paměti.

\begin{table}[htbp]
  \centering
  \caption{Srovnání výkonu analytických dotazů před a po indexaci (průměr z~5 opakování)}
  \label{tab:evaluace:sql}
  \sisetup{group-separator={\,}, output-decimal-marker={,}}

  \begin{tabular}{@{} l p{5cm} S[table-format=5.0] S[table-format=5.0] S[table-format=3.0] S[table-format=2.0] r @{}}
    \toprule
    & & \multicolumn{2}{c}{\textbf{Cena (rel. jednotky)}} & \multicolumn{2}{c}{\textbf{Doba běhu (ms)}} & \\
    \cmidrule(lr){3-4} \cmidrule(lr){5-6}
    \textbf{ID} & \textbf{Popis dotazu} & {\textbf{bez idx}} & {\textbf{s~idx}} & {\textbf{bez idx}} & {\textbf{s~idx}} & \textbf{Zrychl.} \\
    \midrule
    S1 & Počet \PR ve stavu (label) & 22258 & 5424  & 74  & 20 & 3,7$\times$ \\
    S2 & Počet sloučených \PR       & 9598  & 852   & 19  & 3  & 7,0$\times$ \\
    S3 & Soubory upravené týmem     & 16770 & 2653  & 51  & 41 & 1,2$\times$ \\
    S4 & Vývoj počtu \PR v čase     & 24958 & 15120 & 111 & 75 & 1,5$\times$ \\
    \bottomrule
  \end{tabular}
\end{table}